\section{Agente de diseño: del brief a la entrega}

La introducción presentó a Lovart como un Agent vertical; esta sección analiza con más detalle su lógica de producto. El diseño no consiste en generar una imagen una sola vez, sino en un flujo de trabajo de varias rondas: entender el brief, reunir materiales, proponer conceptos, generar propuestas, iterar según la retroalimentación y, al final, entregar archivos utilizables. Lo que el diseñador necesita no es solo «generar una imagen», sino un asistente que entienda el contexto, tenga criterio estético y pueda iterar de forma continua.

\begin{figure}[H]
\centering
\includegraphics[width=0.96\textwidth]{design-agent-workflow.png}
\caption{Flujo de trabajo de un Agent de diseño: del brief a los materiales, las propuestas, la iteración y la entrega. Diagrama conceptual propio, elaborado a partir de la conversación de 00:52:20--00:57:00.}
\end{figure}

\begin{knowledgebox}{Lectura de la figura: por qué el flujo de trabajo de diseño se presta a un Agent vertical}
Las tareas de diseño tienen un contexto definido, exigen un juicio estético fuerte, pasan por varias rondas de retroalimentación y deben producir entregables concretos. Un modelo general puede generar materiales, pero el Agent vertical se encarga de enlazar esos materiales con los objetivos del usuario en un flujo de trabajo continuo.
\end{knowledgebox}

\begin{importantbox}{La naturaleza de un Agent de diseño}
Un Agent de diseño no es un generador de imágenes, sino un coordinador del flujo de trabajo de diseño. Debe reunir en un mismo ciclo los objetivos del usuario, las restricciones de marca, los materiales, el juicio estético y el formato de entrega.
\end{importantbox}

\subsection{Las necesidades AI Native de los diseñadores}

La subsección anterior trató el flujo de trabajo; esta examina en concreto por qué los diseñadores necesitan herramientas AI Native. Los diseñadores no usan la IA solo para ahorrar tiempo, sino también para ampliar su espacio de exploración. Lovart debe ocuparse del contexto de marca, la gestión de materiales, la iteración de versiones, la colaboración y la entrega final. Cuanto más se parezca a una mesa de trabajo profesional, más difícil será que una caja de chat general la sustituya por completo.

\begin{figure}[H]
\centering
\includegraphics[width=0.96\textwidth]{designers-ai-native.png}
\caption{Necesidades AI Native de los diseñadores: los escenarios de diseño profesional requieren criterio estético, contexto, versiones y entrega. Diagrama conceptual propio, elaborado a partir de la conversación de 00:33:00--00:57:00.}
\end{figure}

\begin{knowledgebox}{Términos clave: la pila de capacidades de un Agent de diseño}
\begin{tabularx}{\textwidth}{p{0.22\textwidth}p{0.34\textwidth}X}
\toprule
Capacidad & Función & Dificultad \\
\midrule
Comprensión del brief & Convertir una necesidad expresada en lenguaje natural en objetivos de diseño & Lo que expresa el usuario suele ser ambiguo. \\
Contexto de materiales & Gestionar la marca, las imágenes, las referencias y las restricciones & La información es abundante y cambia constantemente. \\
Juicio estético & Elegir el estilo y la dirección visual & Es muy subjetivo y difícil de evaluar. \\
Iteración de versiones & Modificar y comparar en varias rondas & Hay que recordar la intención de las versiones anteriores. \\
Archivos de entrega & Producir activos de diseño utilizables & El formato, la colaboración y la edición posterior son muy importantes. \\
\bottomrule
\end{tabularx}
\end{knowledgebox}

\begin{knowledgebox}{Las zonas difíciles del diseño}
\begin{tabularx}{\textwidth}{p{0.22\textwidth}p{0.34\textwidth}X}
\toprule
Zona difícil & Manifestación concreta & Por qué no basta un modelo general \\
\midrule
Coherencia de marca & Colores, tipografías, tono y reglas de los activos & Requiere contexto de largo plazo y gestión de activos. \\
Exploración de varias propuestas & Comparar rápidamente direcciones distintas para un mismo brief & Requiere control de versiones y un mecanismo de revisión. \\
Colaboración en equipo & Diseñadores, equipo de operaciones, clientes y jefes modifican juntos & Requiere comentarios, permisos e historial. \\
Formato de entrega & Archivos editables, paquetes de materiales, adaptación de tamaños & Requiere integrarse en la cadena real de herramientas de producción. \\
\bottomrule
\end{tabularx}
\end{knowledgebox}

\begin{importantbox}{La verdadera dificultad del flujo de trabajo de diseño}
La dificultad del trabajo de diseño no está en generar una imagen «bonita», sino en lograr que un conjunto de activos visuales se mantenga coherente frente a la marca, el usuario objetivo, las restricciones del proyecto y la retroalimentación del equipo. Si Lovart quiere convertirse en una mesa de trabajo, tiene que gestionar este contexto continuo.
\end{importantbox}

\subsection{Resumen de la sección}

El valor del producto de Lovart está en llevar la capacidad generativa de la IA al flujo de trabajo de diseño, en lugar de quedarse en la generación de una sola vez. La barrera de entrada de un Agent vertical proviene del contexto profesional y de la entrega continua.
